% Fichier GÉNÉRÉ par scripts/build_report.py depuis report/template/rapport-TP-fr.tex.
% Ne pas éditer : modifier sections/, report/report.toml ou la template, puis `check`.
% Formatage général
	\documentclass[a4paper,11pt]{article}
	\usepackage[T1]{fontenc} % Encodage des caractères spéciaux
	\usepackage{helvet} % Police équivalent Arial
	\usepackage[french]{babel} % Localisation en français
	\frenchbsetup{StandardLists=true} % Pour dire à babel de garder le formatage standard (anglais) des listes à puce
	\usepackage{csquotes} % Pour remplacer automatiquement les guillemets par des guillemets typographique
	\MakeOuterQuote{"} % Le caractère utilisé pour rentrer les guillemets dans le code est le "
	
% Dimensions
	\usepackage[margin=2cm]{geometry} % Marges
	\setlength{\parindent}{0pt} % Largeur de l'alinéa

% Couleurs
	\usepackage{xcolor} % Permet de passer un texte en couleur avec la commande \textcolor{color}{text}
	\definecolor{darkpowderblue}{rgb}{0.0, 0.2, 0.6} % Couleur perso bleue foncée appliquée aux liens URL
	\definecolor{gamboge}{rgb}{0.89, 0.61, 0.06} % Couleur perso orange appliquée aux subsubsections
	\definecolor{ece}{RGB}{0, 122, 123} % Couleur ECE appliquée aux autres titres
	
% Formatage des titres
	\usepackage{titlesec} % Package utilisé pour formater les titres
	\setcounter{secnumdepth}{3} % Numéortation jusqu'au niveau subsubsection
	
	\titleformat{\section} % Formatage des sections
		{\fontsize{18}{18}\bfseries\color{ece}} % Titre en 18 pt, gras, couleur ECE
		{\Roman{section}.} % Numérotation romaine majuscule suivie d'un point
		{1ex} % Un espace entre le numéro et le titre
		{}
	
	\titleformat{\subsection} % Formatage des sous-sections
		{\fontsize{16}{16}\bfseries\color{ece}} % Titre en 16 pt, gras, couleur ECE
		{\Alph{subsection}.} % Numérotation en lettre majuscule suivie d'un point
		{1ex} % Un espace entre le numéro et le titre
		{}
	
	\titleformat{\subsubsection} % Formatage des sous-sous-sections
		{\fontsize{15}{15}\bfseries\color{gamboge}} % Titre en 15 pt, gras, couleur orange
		{(\alph{subsubsection})} % Numérotation en lettre minuscule entourée de parenthèses
		{1ex} % Un espace entre le numéro et le titre
		{}
	
% En-tête et pied de page
	\usepackage{fancyhdr} % Package utilisé pour personnaliser facilement en-tête et pied-de-page
	\renewcommand{\headrulewidth}{0pt} % Épaisseur de la ligne de séparation entre en-tête et corps du rapport
	
	\fancypagestyle{main}{ % En-tête et pied de page pour toutes les autres pages
		\fancyhead[L]{\includegraphics[width=2cm]{images/ece.png}} % En haut à gauche, le logo ECE
		\fancyhead[R]{\textcolor{ece}{CodeTrace} \textbullet{} V1}
		\fancyfoot[C]{\thepage} % En bas au centre, le numéro de page
	}
	\pagestyle{main}
		
% Page de titre
	\newcommand{\HRule}[1]{\rule{\linewidth}{#1}} % Ligne de séparation horizontale d'une épaisseur donnée

% Figures
	\usepackage{graphicx} % Permet l'ajout d'images en spécifiant la taille
	\usepackage{subcaption} % Package utilisé pour avoir des sous-figures avec chacune une légende
	\usepackage[labelfont=bf]{caption} % "Figure" et numéro en gras
	
% Liens
	\usepackage[colorlinks=true,linkcolor=black,urlcolor=darkpowderblue]{hyperref}

% ---------------------------------------------------------------------
% Ajouts recherche (build_report.py).
% Le préambule ci-dessus est celui de la template ECE.
% =====================================================================
% Ajouts "rapport de recherche" au préambule de la template ECE.
% Chargé par main.tex APRÈS le préambule ECE (donc après babel, csquotes, hyperref).
% Ne pas modifier dans report/V*/ : la source est
% .claude/skills/build-report/assets/preamble-recherche.tex
% =====================================================================

% Typographie et maths
	\usepackage{microtype}
	\usepackage{amsmath,amssymb}
	\usepackage[locale=FR]{siunitx} % \num{0.62} -> 0,62 ; \SI, \qty
	\sisetup{group-minimum-digits=4, separate-uncertainty=true}

% Tableaux
	\usepackage{booktabs} % \toprule \midrule \bottomrule, jamais de traits verticaux
	\usepackage{multirow}
	\usepackage{array}
	\usepackage{tabularx}

% Courts extraits de code (uniquement dans des figures, exemples de perturbation)
	\usepackage{listings}
	\lstset{
		basicstyle=\ttfamily\small,
		columns=fullflexible,
		keepspaces=true,
		breaklines=true,
		frame=single,
		rulecolor=\color{gray!50},
		numbers=none,
		showstringspaces=false,
		extendedchars=true,
		inputencoding=utf8,
		literate={é}{{\'e}}1 {è}{{\`e}}1 {à}{{\`a}}1 {ç}{{\c{c}}}1 {ê}{{\^e}}1 {ù}{{\`u}}1
	}

% Bibliographie (biber)
	\usepackage[backend=biber, style=numeric-comp, sorting=nyt, maxbibnames=6, giveninits=true]{biblatex}
	\addbibresource{refs.bib}

% Références croisées : \cref{fig:x} -> "figure 3". Doit venir après hyperref.
	\usepackage[french,noabbrev]{cleveref}
	% babel french rend « : » actif : les labels sec:x cassent \cref (\csname r@sec:x@cref).
	% On neutralise les caractères actifs dans l'argument, comme babel le fait pour \ref (safe=ref).
	\makeatletter
	\AtBeginDocument{%
		\let\ct@cref\cref \let\ct@Cref\Cref
		\renewcommand{\cref}[1]{\begingroup\@safe@activestrue\edef\x{\endgroup\noexpand\ct@cref{#1}}\x}%
		\renewcommand{\Cref}[1]{\begingroup\@safe@activestrue\edef\x{\endgroup\noexpand\ct@Cref{#1}}\x}%
	}
	\makeatother

% Numérotation : les titres affichent \thesection, donc les renvois (\cref) et les titres
% donnent le même numéro (romain pour les parties, lettre pour les annexes).
	\renewcommand{\thesection}{\Roman{section}}
	\renewcommand{\thesubsection}{\thesection.\Alph{subsection}}
	\renewcommand{\thesubsubsection}{\thesubsection.\alph{subsubsection}}
	\titleformat{\section}{\fontsize{18}{18}\bfseries\color{ece}}{\thesection.}{1ex}{}
	\titleformat{\subsection}{\fontsize{16}{16}\bfseries\color{ece}}{\Alph{subsection}.}{1ex}{}
	\titleformat{\subsubsection}{\fontsize{15}{15}\bfseries\color{gamboge}}{(\alph{subsubsection})}{1ex}{}
	\usepackage{appendix}
	\makeatletter
	\g@addto@macro\appendix{%
		\renewcommand{\thesection}{\Alph{section}}%
		\titleformat{\section}{\fontsize{18}{18}\bfseries\color{ece}}{Annexe \thesection.}{1ex}{}%
	}
	\makeatother

% Table des matières
	\setcounter{tocdepth}{2}
	% Numéros romains (VIII, VIII.B) : élargir les cases de numéro pour éviter les overfull.
	\usepackage[titles]{tocloft}
	\setlength{\cftsecnumwidth}{2.6em}
	\setlength{\cftsubsecindent}{2.6em}
	\setlength{\cftsubsecnumwidth}{3.4em}

% Marqueurs de travail (comptés par `build_report.py check`)
	% \afaire{...} : section ou paragraphe prévu mais sans matière réelle dans le repo.
	\newcommand{\afaire}[1]{%
		\par\vspace{2mm}\noindent
		\fcolorbox{gamboge}{gamboge!10}{\parbox{\dimexpr\linewidth-2\fboxsep-2\fboxrule\relax}{%
			\textbf{\textcolor{gamboge}{À venir.}} #1}}%
		\par\vspace{2mm}}
	% \averifier{...} : affirmation dont la source n'est pas encore vérifiée (ex. chiffre lu sur le web).
	\newcommand{\averifier}[1]{#1\textsuperscript{\textcolor{red}{[à vérifier]}}}

% Résultats chiffrés : jamais tapés à la main.
	% \res{run}{cle} affiche la valeur générée par `build_report.py export-results`
	% dans generated/*.tex. Une clé inconnue s'affiche ?? en rouge et lève un warning.
	\makeatletter
	\newcommand{\res}[2]{%
		\ifcsname res@#1@#2\endcsname
			\csname res@#1@#2\endcsname
		\else
			\textbf{\textcolor{red}{??}}%
			\PackageWarning{codetrace}{Resultat inconnu: #1/#2}%
		\fi}
	\newcommand{\defres}[3]{\expandafter\def\csname res@#1@#2\endcsname{#3}}
	\makeatother
	\IfFileExists{generated/resultats.tex}{\input{generated/resultats.tex}}{}

% Figures et tableaux : légende au-dessus des tableaux, en dessous des figures
	\captionsetup[table]{position=top}

\hypersetup{pdftitle={CodeTrace : Détection de la provenance des données d'entraînement des LLM pour le code}, pdfauthor={À compléter}, pdflang={fr}}

\begin{document}

% Page de titre, reprise de la mise en page ECE (logo, filets, couleur ECE).
% Généré par `build_report.py new` depuis report/report.toml : ne pas éditer à la main.
\begin{titlepage}
	\begin{minipage}[t]{.45\textwidth}
		\includegraphics[width=6cm]{images/ece.png}
	\end{minipage}
	\hfill
	\begin{minipage}[t]{.45\textwidth}
		\raggedleft\vspace{-12mm}\Large{\textbf{Projet de recherche}\\ 2026-2027}
	\end{minipage}

	\vspace{3cm}
	\centering
	\HRule{1.5pt} \\ [0.5cm]
	{\Large\textbf{\MakeUppercase{Rapport de recherche}}}\\ [5mm]
	{\huge\textbf{\textcolor{ece}{CodeTrace}}}\\ [4mm]
	{\LARGE\textcolor{ece}{Détection de la provenance des données d'entraînement des LLM pour le code}}\\ [3mm]
	\HRule{1.5pt} \\ [1cm]

	{\Large Version 1}\\ [2mm]
	{\large 14 novembre 2026}

	\vspace{2cm}
	{\large\textbf{Auteurs}}\\ [2mm]
	{\large À compléter}

	\vspace{1cm}
	{\large\textbf{Encadrants}}\\ [2mm]
	{\large Subhankar Maity \\ Guilherme Medeiros Machado}

	\vfill
	Équipe Data \& IA, ECE, Paris
\end{titlepage}


% Résumé (FR) et abstract (EN). Réécrits à CHAQUE version pour refléter l'état réel du projet.
% 150 à 250 mots chacun : contexte, question, méthode, résultat principal (ou état d'avancement), portée.
% Sources : tout le rapport (V1 : CLAUDE.md, notes/etat_de_l_art.md, notes/journal.md).
\section*{Résumé}
\addcontentsline{toc}{section}{Résumé}
Les grands modèles de langage sont entraînés sur de vastes corpus de code public, souvent sans que les auteurs de ce code le sachent ni y consentent. Savoir si un extrait de code donné a servi à l'entraînement d'un modèle intéresse la protection des licences, la détection de fuites de secrets, la contamination des bancs d'essai et l'évaluation pédagogique. Le projet CodeTrace pose la question suivante : peut-on déterminer si un extrait de code a fait partie des données d'entraînement d'un modèle en comparant le comportement de ce modèle (log-probabilités des tokens, sorties générées) sur la version authentique et sur des versions perturbées du même code ?

Ce premier rapport décrit l'état du projet en phase de cadrage. Il présente une synthèse préliminaire de la littérature sur l'inférence d'appartenance appliquée aux modèles de langage et au code, fondée sur les résumés des travaux et non encore sur des fiches de lecture détaillées. Il fixe les notations, le cadre formel et les règles méthodologiques retenues : labels membre et non-membre documentés, baseline aveugle contre le biais temporel, métriques avec intervalles de confiance. Aucune donnée n'a encore été collectée, aucune méthode n'est implémentée et aucun résultat expérimental n'est disponible. Les prochaines étapes sont la lecture détaillée des papiers de référence, la vérification de l'accès aux log-probabilités pour chaque modèle visé et le choix de la source des labels.

\noindent\textbf{Mots-clés :} inférence d'appartenance, données d'entraînement, grands modèles de langage, code source, provenance.

\section*{Abstract}
\addcontentsline{toc}{section}{Abstract}
Large language models are trained on vast corpora of public code, often without the knowledge or consent of its authors. Knowing whether a given code snippet was used to train a model matters for licence enforcement, secret leakage detection, benchmark contamination and educational assessment. The CodeTrace project asks the following question: can we determine whether a code snippet was part of a model's training data by comparing the model's behaviour (token log-probabilities, generated outputs) on the authentic version and on perturbed versions of the same code?

This first report describes the project at the scoping stage. It presents a preliminary synthesis of the literature on membership inference for language models and for code, based on the papers' abstracts rather than on detailed reading notes. It sets the notation, the formal framework and the methodological rules adopted: documented member and non-member labels, a blind baseline against temporal bias, and metrics with confidence intervals. No data has been collected yet, no method is implemented and no experimental result is available. The next steps are a detailed reading of the reference papers, checking log-probability access for each target model, and choosing the source of membership labels.

\noindent\textbf{Keywords:} membership inference, training data, large language models, source code, provenance.
\newpage


\tableofcontents
\newpage

% Sources : CLAUDE.md (question de recherche, objectifs), notes/journal.md, notes/etat_de_l_art.md (motivations).
\section{Introduction}\label{sec:introduction}
	\subsection{Contexte et motivations}\label{sec:intro-contexte}
		Les grands modèles de langage (\emph{large language models}, LLM) sont entraînés sur des corpus qui contiennent une grande quantité de code source public. La composition exacte de ces corpus est rarement publiée. Un développeur ne peut donc pas savoir directement si son code a servi à entraîner un modèle donné.

		Cette question a plusieurs enjeux. Le premier est juridique : du code publié sous licence (GPL, CC-BY-SA) peut avoir été utilisé sans respect des conditions de la licence. Le deuxième concerne la sécurité : un dépôt qui a contenu des secrets (clés d'API, identifiants) a pu exposer ces secrets au modèle. Le troisième est scientifique : un banc d'essai de génération de code dont les solutions figurent dans les données d'entraînement ne mesure plus la capacité de généralisation du modèle (contamination). Le quatrième est pédagogique : évaluer un exercice de programmation suppose de savoir si sa solution est déjà connue du modèle.

		Le problème relève de l'inférence d'appartenance (\emph{membership inference attack}, MIA) : décider, à partir du comportement d'un modèle, si un échantillon faisait partie de ses données d'entraînement. Le code présente des particularités qui le distinguent du langage naturel (syntaxe contraignante, forte duplication, transformations qui préservent exactement la sémantique), présentées en \cref{sec:ea-code}.

	\subsection{Question de recherche}\label{sec:intro-question}
		La question de recherche de CodeTrace est la suivante.

		\begin{quote}
			Peut-on déterminer si un extrait de code donné a fait partie des données d'entraînement d'un LLM, en comparant le comportement du modèle (log-probabilités des tokens, sorties) sur la version authentique et sur des versions perturbées du même code ?
		\end{quote}

		L'hypothèse sous-jacente, que nous prévoyons de tester, est qu'un modèle qui a vu un extrait pendant son entraînement le traite de façon plus "familière" que des versions sémantiquement équivalentes mais textuellement différentes de cet extrait. L'écart de comportement entre version authentique et versions perturbées serait alors plus grand pour les membres que pour les non-membres.

		Le sujet fixe sept objectifs : (1) curer des jeux de données de code réels et générer des variantes perturbées ; (2) tester l'appartenance en comparant version authentique et versions perturbées ; (3) évaluer plusieurs LLM (le sujet cite GPT-5, Gemini 2.5 Flash, Mistral 24B, LLaMA 3.3 70B et DeepSeek) ; (4) définir des critères de décision robustes, par modèle et par type de code ; (5) visualiser et comparer les méthodes avec des métriques standard ; (6) construire un outil interactif de test ; (7) analyser l'influence de l'architecture des modèles et des caractéristiques du code.

	\subsection{Contributions}\label{sec:intro-contributions}
		À ce stade du projet, les contributions réalisées sont limitées au cadrage :
		\begin{itemize}
			\item une synthèse préliminaire de la littérature sur les MIA appliquées aux LLM et au code (\cref{sec:etat-art}), rédigée à partir des résumés des travaux ;
			\item un cadre formel et des règles méthodologiques fixées avant toute expérience (\cref{sec:cadre,sec:protocole}), en particulier l'obligation d'une baseline aveugle contre le biais temporel ;
			\item une infrastructure de travail reproductible : dépôt structuré, liste versionnée des papiers de référence avec script de téléchargement, et génération versionnée de ce rapport.
		\end{itemize}
		Aucune donnée, méthode ou résultat expérimental n'est encore disponible.

	\subsection{Organisation du rapport}\label{sec:intro-plan}
		La \cref{sec:etat-art} présente l'état de l'art. La \cref{sec:cadre} formalise le problème et le modèle de menace. Les \cref{sec:donnees,sec:methodes,sec:protocole} décriront les données, les méthodes et le protocole expérimental. Les \cref{sec:resultats,sec:discussion} présenteront et discuteront les résultats, la \cref{sec:outil} l'outil interactif. La \cref{sec:conclusion} conclut. Les annexes consignent les décisions méthodologiques (\cref{ann:decisions}), les informations de reproductibilité (\cref{ann:repro}) et l'historique des versions (\cref{ann:historique}).

% Sources : notes/etat_de_l_art.md (synthèse à partir des résumés, 2026-11-12), papers/papers.toml.
% Aucune fiche de lecture n'existe encore dans notes/ : aucun chiffre de la littérature n'est cité,
% et les conclusions des papiers sont marquées \averifier tant qu'elles ne sont pas fichées.
% Règle : un chiffre d'un papier n'est écrit que si une fiche existe dans notes/ ; sinon \averifier{}.
\section{État de l'art}\label{sec:etat-art}
	Cette section s'appuie sur une synthèse préliminaire rédigée à partir des résumés des travaux. Aucun papier n'a encore fait l'objet d'une fiche de lecture détaillée. Nous décrivons donc les méthodes proposées, mais les conclusions empiriques des auteurs sont signalées comme restant à vérifier, et aucun chiffre de la littérature n'est reproduit.

	Nous distinguons trois niveaux de granularité. Au niveau de l'échantillon, on décide si un extrait donné est membre des données d'entraînement : c'est le cadre classique des MIA, aussi appelé détection de données de pré-entraînement (\emph{pretraining data detection}). Au niveau du jeu de données (\emph{dataset inference}), on décide si un ensemble d'extraits (par exemple un dépôt) a été utilisé. Enfin, le traçage actif consiste à marquer son code avant publication pour pouvoir prouver plus tard son utilisation.

	\subsection{Inférence d'appartenance sur les grands modèles de langage}\label{sec:ea-mia}
		Les premières attaques exploitent la perte du modèle. Un membre tend à avoir une perte (ou une perplexité) plus faible qu'un non-membre. Carlini et al.~\cite{carlini2021} utilisent ce signal pour extraire des données d'entraînement et le normalisent, entre autres, par l'entropie de compression zlib du texte : un texte très compressible est facile à prédire pour tout modèle, ce qui ne prouve pas qu'il a été vu.

		Une autre famille compare la perte du modèle cible à celle d'un modèle de référence entraîné sur des données de distribution proche. Cette calibration est coûteuse et suppose de disposer de données de même distribution, ce qui est souvent irréaliste. Mattern et al.~\cite{mattern2023} proposent de remplacer le modèle de référence par des voisins de l'échantillon : ils génèrent des versions légèrement perturbées de $x$ (remplacement de mots) et comparent la perte sur $x$ à la perte moyenne sur ces voisins. Cette comparaison avec des voisins perturbés est l'idée la plus proche de celle de CodeTrace.

		Shi et al.~\cite{shi2024} proposent Min-K\% Prob, fondé sur l'intuition qu'un texte non vu contient quelques tokens aberrants de faible probabilité : le score est la moyenne des $k\,\%$ log-probabilités les plus faibles de l'extrait. Le même travail introduit le banc d'essai WikiMIA, dont les non-membres sont définis par un critère temporel. Zhang et al.~\cite{zhang2025} proposent Min-K\%++, qui normalise la log-probabilité de chaque token par la moyenne et l'écart-type de la distribution conditionnelle du modèle sur le vocabulaire, en s'appuyant sur l'idée que les échantillons d'entraînement tendent à être des maxima locaux de la distribution apprise. \averifier{Les auteurs présentent cette méthode comme l'état de l'art sans modèle de référence sur WikiMIA et MIMIR.}

		Ces méthodes ont une exigence commune : elles demandent les log-probabilités des tokens de l'extrait fourni en entrée, et même la distribution complète sur le vocabulaire pour Min-K\%++. Cette exigence conditionne les modèles auxquels elles s'appliquent (\cref{sec:ea-boite-noire}).

	\subsection{Limites méthodologiques des évaluations}\label{sec:ea-critique}
		Plusieurs travaux récents remettent en cause la validité des évaluations de MIA sur les LLM. Ce point conditionne directement notre protocole.

		Duan et al.~\cite{duan2024} évaluent les MIA sur des membres et non-membres tirés de la même distribution (banc d'essai MIMIR, modèles Pythia entraînés sur la Pile). \averifier{Ils concluent que les MIA y font à peine mieux que le hasard, et attribuent une partie des succès rapportés ailleurs à un décalage de distribution, notamment temporel, entre membres et non-membres.} Das et al.~\cite{das2024} construisent des classifieurs \emph{aveugles}, qui n'interrogent pas le modèle. \averifier{Ces baselines aveugles battent les meilleures MIA publiées sur plusieurs bancs d'essai, ce qui indique que ces bancs d'essai séparent les classes par leur distribution plutôt que par l'appartenance.} Meeus et al.~\cite{meeus2024} dressent un bilan allant dans le même sens et proposent des protocoles plus sûrs, comme des partitions aléatoires contrôlées ou des données dont la date d'inclusion est certaine.

		Maini et al.~\cite{maini2024} proposent de passer au niveau du jeu de données : plusieurs scores MIA sont agrégés, puis un test statistique est appliqué à un ensemble suspect. Ce cadre correspond bien au cas d'usage "mon dépôt a-t-il servi à l'entraînement ?".

		Nous en tirons deux règles pour CodeTrace. Une baseline aveugle, qui distingue code ancien et code récent sans interroger le modèle, est obligatoire. Les labels membre et non-membre doivent, autant que possible, provenir d'un corpus d'entraînement connu plutôt que de la seule date de publication (\cref{sec:donnees-labels}).

	\subsection{Travaux spécifiques au code source}\label{sec:ea-code}
		Al-Kaswan et al.~\cite{alkaswan2024} étudient la mémorisation dans les LLM de code par des attaques d'extraction. \averifier{Ils rapportent que la mémorisation augmente avec la taille du modèle et que les fichiers de données et constantes (\emph{data carriers}) sont plus mémorisés que le code ordinaire.}

		Wan et al.~\cite{wan2024} adaptent les MIA aux modèles de complétion de code en boîte noire, à l'aide de modèles fantômes (\emph{shadow models}). \averifier{L'attaque fonctionne sur des modèles de petite taille mais l'appartenance reste difficile à détecter sur des modèles plus grands comme CodeGen et StarCoder.} Yang et al.~\cite{yang2024} et Zhang et Li~\cite{zhang2023} proposent d'autres MIA pour les modèles de code, respectivement avec un modèle substitut et en boîte noire.

		Jiang et al.~\cite{jiang2025} sont les plus proches de l'approche "authentique contre perturbé" de CodeTrace. Ils génèrent des perturbations du code qui préservent la sémantique (prompts adverses), comparent les complétions obtenues au code réel et entraînent un classifieur sur les caractéristiques qui en résultent. Leur méthode reste centrée sur des modèles de complétion de code.

		Deux travaux trouvés sur le web, non encore lus, exploitent la structure du code. Li et al.~\cite{li2025} retirent avant le calcul du score les tokens imposés par la syntaxe de Python, prévisibles pour tout modèle, et proposent un banc d'essai Python. Yang et al.~\cite{yang2025} étudient l'effet de transformations sémantiquement équivalentes (renommage, reformatage) sur les MIA. \averifier{Le renommage de variables y ferait baisser le succès des attaques}, ce qui montrerait que la perturbation peut servir à la fois d'outil de détection et de moyen d'évasion.

		Le code se distingue du langage naturel sur trois points. Une grande part des tokens est déterminée par la syntaxe. Le code est très dupliqué (forks, code générique, en-têtes de licence), ce qui brouille la frontière entre membre et non-membre. Enfin, des transformations préservent exactement la sémantique (renommage, reformatage, réordonnancement, commentaires) : elles fournissent des voisins au sens de Mattern et al.~\cite{mattern2023} sans recourir à un modèle de langue pour les générer.

	\subsection{Accès boîte noire et API fermées}\label{sec:ea-boite-noire}
		Les scores des \cref{sec:ea-mia,sec:ea-code} supposent l'accès aux log-probabilités du code fourni en entrée. \averifier{Les API commerciales (OpenAI, Gemini) semblent n'exposer que les log-probabilités des tokens générés, en nombre limité.} Ce point sera vérifié modèle par modèle (\cref{sec:cadre-acces}).

		Plusieurs travaux récents ciblent le cadre où l'on n'observe que du texte. Hallinan et al.~\cite{hallinan2025} donnent un préfixe au modèle, échantillonnent des continuations et mesurent leur recouvrement en n-grammes avec le vrai suffixe. Yi et Li~\cite{yi2026} estiment le comportement local du modèle à partir du seul mot suivant échantillonné. He et al.~\cite{he2025} proposent une attaque qui n'utilise que les tokens générés (\emph{label-only}), fondée sur leur similarité sémantique. Du côté de l'audit d'API, Chauvin et al.~\cite{chauvin2025} utilisent les log-probabilités renvoyées par une API pour détecter un changement du modèle servi. Un article de l'équipe encadrante sur l'audit d'API de LLM boîte noire par les probabilités des tokens, appliqué à la génération de code, se situe à ce croisement ; sa référence exacte reste à obtenir.

		Pour CodeTrace, cela suggère deux protocoles. Pour les modèles à poids ouverts (\emph{open-weights}), les scores fondés sur les log-probabilités du prompt sont utilisables. Pour les modèles fermés, il faudrait des scores fondés sur la complétion : donner un préfixe du code authentique ou perturbé, puis mesurer la fidélité de la suite générée.

	\subsection{Approches proactives : tatouage du code}\label{sec:ea-tatouage}
		Au lieu d'inférer l'appartenance après coup, on peut marquer le code avant sa publication. CoProtector insère du code mort servant de marque \averifier{(Sun et al., 2022)}. CodeMark~\cite{sun2023} insère des marques imperceptibles par des transformations qui préservent la sémantique. \averifier{L'efficacité de ces marques est contestée par des travaux ultérieurs, et des attaques comme DeCoMa les détectent et les retirent.} Surtout, ces méthodes ne s'appliquent pas au code déjà publié, qui est précisément la cible de CodeTrace.

	\subsection{Positionnement de CodeTrace}\label{sec:ea-positionnement}
		D'après cette synthèse préliminaire, les MIA au niveau de l'échantillon sur de grands LLM pré-entraînés donneraient un signal faible lorsque l'évaluation est rigoureuse, et une partie des résultats positifs s'expliquerait par un biais de distribution. La littérature récente évolue dans trois directions : le niveau du jeu de données, les méthodes sans log-probabilités pour les API fermées, et l'adaptation des scores à la structure du code.

		CodeTrace vise la lacune suivante :
		\begin{itemize}
			\item comparer systématiquement le comportement d'un modèle sur du code authentique et sur des variantes sémantiquement équivalentes, sur des LLM généralistes récents, à poids ouverts et fermés, et non seulement sur des modèles de complétion de code ;
			\item le faire avec des labels fiables, issus de corpus connus, et contre une baseline aveugle ;
			\item analyser l'effet du type de code (langage, code générique, fichiers de données) et de l'architecture du modèle ;
			\item proposer un outil interactif de test.
		\end{itemize}
		Ce positionnement sera révisé après la lecture détaillée des papiers de référence.

% Sources : CLAUDE.md (question de recherche, contraintes d'accès), docs/model_access.md (gabarit vide au 2026-11-14), notes/decisions.md.
\section{Cadre formel et modèle de menace}\label{sec:cadre}
	\subsection{Définition du problème}\label{sec:cadre-definition}
		Soit $M$ un modèle de langage entraîné sur un corpus $D$ inconnu de l'auditeur, et $x$ un extrait de code. La question d'appartenance est de décider si $x \in D$. Une attaque associe à $x$ un score $s(x, M)$, puis décide que $x$ est membre si le score dépasse un seuil :
		\begin{equation*}
			\hat{y}(x) = \mathbb{1}\left[ s(x, M) > \tau \right].
		\end{equation*}
		Par convention, un score plus élevé indique une appartenance plus probable.

		CodeTrace s'intéresse aux scores qui comparent l'extrait authentique $x$ à un ensemble de versions perturbées $\tilde{x}_1, \dots, \tilde{x}_m$, sémantiquement équivalentes à $x$ mais textuellement différentes. Ces scores prennent la forme générale
		\begin{equation*}
			s(x, M) = g\bigl(f(x, M), \{ f(\tilde{x}_j, M) \}_{j=1}^{m}\bigr),
		\end{equation*}
		où $f$ mesure la familiarité du modèle avec un texte (par exemple l'opposé de sa perte) et $g$ compare la valeur sur l'original aux valeurs sur les versions perturbées. Les choix de $f$, de $g$ et des perturbations seront précisés en \cref{sec:methodes}.

		Au niveau du jeu de données, la question porte sur un ensemble $X = \{x_1, \dots, x_n\}$ (par exemple les fichiers d'un dépôt). La décision repose alors sur un test statistique appliqué à la distribution des scores de $X$, comparée à celle d'un ensemble de non-membres connus.

		Le seuil $\tau$ n'est pas universel : il dépend du modèle et potentiellement du type de code. Nous évaluerons donc les scores indépendamment du seuil (AUC-ROC) et à des taux de faux positifs fixés (\cref{sec:protocole-metriques}).

	\subsection{Modèle de menace et accès aux modèles}\label{sec:cadre-acces}
		L'auditeur ne connaît pas $D$ et n'a aucun contrôle sur l'entraînement. Il dispose de l'extrait $x$ et d'un accès en lecture au modèle, dont le niveau varie :
		\begin{itemize}
			\item \textbf{boîte blanche} : accès aux poids et donc aux logits de tous les tokens, en entrée comme en sortie (modèles à poids ouverts exécutés localement) ;
			\item \textbf{boîte grise} : accès par API aux log-probabilités, éventuellement limitées aux tokens générés et aux $k$ tokens les plus probables ;
			\item \textbf{boîte noire} : accès au seul texte généré.
		\end{itemize}

		La distinction critique est celle entre log-probabilités du prompt et log-probabilités de la complétion. Les scores de la \cref{sec:ea-mia} demandent les premières. Si une API ne renvoie que les secondes, il faut une méthode fondée sur la complétion du code.

		\afaire{Tableau des modèles visés (GPT-5, Gemini 2.5 Flash, Mistral 24B, LLaMA 3.3 70B, DeepSeek) avec, pour chacun : mode d'accès, fournisseur, log-probabilités du prompt et de la complétion, top-$k$ maximal, coût, date de cutoff annoncée, source et date de vérification. Rien n'est encore vérifié dans docs/model\_access.md.}

% Sources : data/README.md (aucun dataset au 2026-11-14), CLAUDE.md (règles méthodologiques), notes/decisions.md.
% Toujours documenter comment le label membre / non-membre est obtenu.
\section{Données}\label{sec:donnees}
	Aucun jeu de données n'a encore été collecté. Cette section fixe les exigences que les données devront satisfaire.

	\subsection{Sources de code}\label{sec:donnees-sources}
		\afaire{Corpus utilisés, licences, langages, période couverte, script de reconstruction.}
	\subsection{Obtention des labels membre et non-membre}\label{sec:donnees-labels}
		Pour chaque jeu de données, la manière d'obtenir le label sera documentée explicitement. Trois sources sont envisagées : la date de cutoff annoncée du modèle comparée à la date du commit, l'appartenance à un corpus public connu pour avoir servi à l'entraînement (comme The Stack ou la Pile), ou des données dont la date d'inclusion est certaine. Compte tenu des limites rappelées en \cref{sec:ea-critique}, les corpus connus seront préférés à la seule date.

		\afaire{Choix de la source des labels (décision à consigner dans notes/decisions.md).}
	\subsection{Nettoyage et déduplication}\label{sec:donnees-dedup}
		\afaire{Filtrage, déduplication exacte et approchée entre membres et non-membres.}
	\subsection{Contrôle du biais temporel}\label{sec:donnees-biais}
		Un classifieur qui distingue code ancien et code récent sans interroger le modèle sera entraîné sur les mêmes données que les attaques. Si ce classifieur aveugle atteint une performance comparable à celle des MIA, les données ne permettent pas de conclure sur l'appartenance.

		\afaire{Mise en œuvre et résultats de la baseline aveugle.}

% Sources : src/codetrace/perturb/, src/codetrace/attacks/, notes/decisions.md.
\section{Méthodes}\label{sec:methodes}
	\subsection{Perturbations du code}\label{sec:methodes-perturbations}
		\afaire{Renommage, reformatage, réordonnancement, commentaires. Garanties de préservation sémantique. Exemples en annexe.}
	\subsection{Scores d'appartenance}\label{sec:methodes-scores}
		\afaire{Définitions mathématiques : loss, zlib, Min-K\%, Min-K\%++, comparaison authentique contre perturbé.}
	\subsection{Critères de décision}\label{sec:methodes-decision}
		\afaire{Choix des seuils par modèle et par type de code.}

% Sources : CLAUDE.md (règles méthodologiques, contraintes d'accès), docs/model_access.md (non rempli au 2026-11-14).
% À venir : experiments/*.yaml, src/codetrace/eval/.
\section{Protocole expérimental}\label{sec:protocole}
	Aucune expérience n'a encore été lancée. Les règles ci-dessous sont fixées avant toute expérience.

	\subsection{Modèles évalués}\label{sec:protocole-modeles}
		Le sujet cite GPT-5, Gemini 2.5 Flash, Mistral 24B, LLaMA 3.3 70B et DeepSeek.

		\afaire{Versions exactes, mode d'accès (local ou fournisseur), date des appels, après vérification de l'accès aux log-probabilités (\cref{sec:cadre-acces}).}
	\subsection{Métriques et tests statistiques}\label{sec:protocole-metriques}
		Chaque score sera évalué par l'aire sous la courbe ROC (AUC-ROC) et par le taux de vrais positifs (TPR) à 1\,\% et 5\,\% de faux positifs (FPR). Le régime à faible taux de faux positifs est le plus pertinent pour un audit, où une accusation erronée a un coût élevé. Chaque métrique sera accompagnée d'un intervalle de confiance obtenu par bootstrap. La comparaison de deux méthodes sur les mêmes extraits reposera sur un test explicite (test de Wilcoxon apparié ou test de permutation).
	\subsection{Baselines}\label{sec:protocole-baselines}
		La baseline aveugle décrite en \cref{sec:donnees-biais} est obligatoire. Les scores de la littérature (perte, ratio zlib, Min-K\%, Min-K\%++) serviront de points de comparaison pour les scores fondés sur les perturbations.

		\afaire{Implémentation des baselines.}
	\subsection{Reproductibilité et coût}\label{sec:protocole-repro}
		Les seeds seront fixées et les configurations versionnées. Chaque expérience produira un dossier daté dans \texttt{results/} contenant une copie de sa configuration. Toutes les réponses d'API seront mises en cache sur disque, avec une clé formée du modèle et d'un hachage du prompt, afin de ne jamais payer deux fois la même requête. Le coût de chaque campagne d'appels sera estimé avant son lancement. Aucun chiffre de ce rapport ne sera recopié à la main : les résultats seront exportés automatiquement depuis \texttt{results/}.

		\afaire{Coût estimé et réel de chaque campagne.}

% Sources : results/<date>_<nom>/ UNIQUEMENT, via generated/ (`build_report.py export-results`).
% Aucun nombre tapé à la main : utiliser \res{run}{cle} et \input{generated/...}.
% Passer les nouveaux résultats au subagent stats-reviewer avant de les écrire ici.
\section{Résultats}\label{sec:resultats}
	\afaire{Aucun résultat expérimental pour l'instant.}

% Sources : résultats, retours du subagent stats-reviewer, notes/decisions.md.
\section{Discussion et limites}\label{sec:discussion}
	\subsection{Interprétation}\label{sec:discussion-interpretation}
		\afaire{Interprétation des résultats au regard de la question de recherche.}
	\subsection{Menaces à la validité}\label{sec:discussion-validite}
		\afaire{Validité interne (labels, biais temporel, fuite), externe (modèles, langages), de construction (métriques).}

% Sources : app/.
\section{Outil interactif}\label{sec:outil}
	\afaire{Architecture, fonctionnement, captures d'écran.}

% Sources : synthèse du rapport, notes/journal.md (prochaines étapes du 2026-11-12).
\section{Conclusion et perspectives}\label{sec:conclusion}
	Le projet CodeTrace est en phase de cadrage. La question de recherche, le cadre formel et les règles méthodologiques sont fixés. Une synthèse préliminaire de la littérature suggère que les MIA au niveau de l'échantillon sur de grands LLM donnent un signal faible lorsque l'évaluation contrôle le biais de distribution ; cette conclusion reste à confirmer par la lecture détaillée des papiers. Aucune donnée, méthode ni résultat n'est encore disponible.

	Les prochaines étapes, par ordre de priorité, sont les suivantes :
	\begin{enumerate}
		\item rédiger une fiche de lecture pour chacun des huit papiers de référence, ce qui permettra de lever les marqueurs "à vérifier" de la \cref{sec:etat-art} ;
		\item vérifier pour chaque modèle visé si l'API renvoie les log-probabilités du prompt ou seulement celles de la complétion, et à quel coût (\cref{sec:cadre-acces}) ;
		\item choisir la source des labels membre et non-membre pour le code, entre dates de cutoff des modèles et de commits, et corpus d'entraînement connus (\cref{sec:donnees-labels}).
	\end{enumerate}


\newpage
\printbibliography[heading=bibintoc, title={Références}]

\newpage
\appendix
% Source : notes/decisions.md (une entrée par décision : date, décision, justification, alternatives écartées).
\section{Décisions méthodologiques}\label{ann:decisions}
	\afaire{Aucune décision consignée dans notes/decisions.md.}

% Sources : version.json, experiments/, results/*/config copiée.
\section{Reproductibilité}\label{ann:repro}
	Cette version du rapport a été générée à partir du commit \texttt{befd249} du dépôt CodeTrace. Elle ne cite aucun résultat expérimental.

	\afaire{Pour chaque résultat : run, config, commit, seed.}

% Source : CHANGELOG.md de chaque version (la plus récente en premier).
% Une sous-section par version : date, commit de base, principales modifications.
\section{Historique des versions}\label{ann:historique}
	\subsection{Version 1 (14 novembre 2026)}
		Première version, générée depuis le commit \texttt{befd249}. Elle contient le résumé, l'introduction, une synthèse préliminaire de l'état de l'art fondée sur les résumés des travaux, le cadre formel, les exigences sur les données et les règles du protocole expérimental. Les sections consacrées aux méthodes, aux résultats, à la discussion et à l'outil interactif sont des emplacements réservés.

% Sources : src/codetrace/perturb/, tests. Courts exemples authentique / perturbé en figures.
\section{Exemples de perturbations}\label{ann:perturbations}
	\afaire{Exemples authentique et perturbé côte à côte.}


\end{document}
